
本书用各类伪全纯圆盘的模空间及其基本链定义各种算子、同态等等。
为了定义这些基本链，我们需要适当地扰动模空间，使其具有“正确维数”
（即伪全纯曲线方程线性化算子的指标）。

先验地说，我们的模空间未必横截。我们采用 
\textbf{[FuOn99II]} 所发展的普遍框架（相关工作见
\textbf{[LiTi98, Rua99, Sie96]}）；我们在那里引入了仓西结构的概念。
带仓西结构空间的典型例子是各类模空间；在这种情形下，仓西邻域由
定义模空间的微分方程作有限维约化而构造。不过，仓西结构的定义还适用于
研究模空间以外更一般的情形。

在仓西结构的框架中，我们把障碍丛的多值截面用作原方程的多值扰动，
由此构造具有正确维数的扰动后模空间及该空间的基本链。我们把这一基本链
称为虚基本链。这套机制适用于几何中出现的许多不同情形。本附录相当完整
而自足地说明本书应用所需的仓西结构理论。除若干技术性引理的证明须参见
原论文 \textbf{[FuOn99II]} 外，我们给出仓西结构各项基本性质的大部分证明。
仓西结构及其多值截面是颇为初等的概念；理解它们只需要拓扑与几何的
基础知识。

本附录的大部分内容是对 \textbf{[FuOn99II]} 材料的回顾。因此，熟悉
仓西结构的读者不必阅读第~\ref{sec:a1-1} 节。不过，我们在 A1.2、A1.3
两节中为 \textbf{[FuOn99II]} 已给出的内容增加了几点：A1.2 节定义带仓西
结构空间的纤维积，A1.3 节定义仓西结构上的有限群作用。

我们归纳构造的各种模空间之边界，都是先前归纳步骤所构造模空间的纤维积
（见第~\ref{sec:7-1} 节）。因此，要精确描述模空间的边界，纤维积起着
本质作用。我们还指出，在 \textbf{[FuOn99II]} 对一般辛流形构造
格罗莫夫--威腾不变量时，模空间的紧化也以类似的归纳方式使用纤维积。
不过，在格罗莫夫--威腾不变量的情形中，可以在同调层次工作，从而可在
每一阶段分别选取扰动（即多值截面）；所以纤维积在紧化中的使用并非如此
本质。另一方面，本书必须在链层次工作，关键之处就在于使不同归纳步骤中
构造的多值截面系彼此相容。因此，必须精确定义这种相容性；为此又必须
结合横截性与胶合，细致处理模空间的纤维积。

下面简述本附录的内容。A1.1 节快速回顾 \textbf{[FuOn99II]} 中关于仓西
结构的材料；A1.2 节定义纤维积，A1.3 节定义仓西结构上的有限群作用。
其实，\textbf{[FuOn99II]} 第 5.4 节已经使用过后一概念，不过这里会给出
更精确的描述。A1.4 节补充讨论仓西结构坐标变换的光滑性。A1.5 节考察
若干（反）例，用以说明仓西结构与轨形定义中出现的一些微妙问题。

\textbf{[FOOO00]} 中有若干技术性错误，\textbf{[FuOn99II]} 中也有一处
与本附录材料有关的错误。本版已经改正所有这些错误，而且没有一个是本质
性的。为方便已经读过 \textbf{[FOOO00, FuOn99II]} 的读者，我们在 A1.6
节中指出这些错误及其改正办法；其他读者无须阅读该节。

\section{仓西结构与多值截面定义回顾}
\label{sec:a1-1}

首先定义（带角的）仓西结构。这里不用“轨形的光滑嵌入”这一说法，
重新表述 \textbf{[FuOn99II]} 中的定义。设 $X$ 是紧可度量空间，且
$p\in X$。

\begin{definition}{A1.1}\label{def:a1-1}
$X$ 中 $p$ 的一个\emph{仓西邻域}\liftIndex{I308}是五元组
$(V_p,E_p,\Gamma_p,\psi_p,s_p)$，满足：
\begin{description}[style=multiline,leftmargin=7.2em,labelwidth=6.5em]
\item[\textup{(A1.2.1)}] $V_p$ 是有限维光滑流形，可以有边界或角，也可以没有；
\item[\textup{(A1.2.2)}] $E_p$ 是有限维实向量空间；
\item[\textup{(A1.2.3)}] $\Gamma_p$ 是在 $V_p$ 上光滑、有效作用的有限群，并在
$E_p$ 上有一个线性表示；
\item[\textup{(A1.2.4)}] $s_p\colon V_p\to E_p$ 是 $\Gamma_p$-等变光滑映射；
\item[\textup{(A1.2.5)}] $\psi_p$ 是从 $s_p^{-1}(0)/\Gamma_p$ 到 $X$ 中 $p$
的一个邻域上的同胚。
\end{description}
\end{definition}

置 $U_p=V_p/\Gamma_p$，并称 $U_p$ 为一个仓西邻域。滥用记号时，
我们有时也称 $V_p$ 为仓西邻域。我们把
$E_p\times V_p\to V_p$ 称为\emph{障碍丛}\liftIndex{I310}%
\liftIndex{I316}，把 $s_p$ 称为\emph{仓西映射}\liftIndex{I210}。
对 $x\in V_p$，以 $(\Gamma_p)_x$ 表示 $x$ 处的稳定子群，即
\[
 (\Gamma_p)_x=\{\gamma\in\Gamma_p\mid \gamma x=x\}.
\]

取 $o_p\in V_p$，使 $s_p(o_p)=0$ 且 $\psi_p([o_p])=p$。我们可以并且
将会假定 $o_p$ 被 $\Gamma_p$ 的每个元素固定。

\begin{definition}{A1.3}\label{def:a1-3}
设 $(V_p,E_p,\Gamma_p,\psi_p,s_p)$ 与
$(V_q,E_q,\Gamma_q,\psi_q,s_q)$ 分别是 $p\in X$ 与
$q\in\psi_p(s_p^{-1}(0)/\Gamma_p)$ 的仓西邻域。若三元组
$(\widehat\phi_{pq},\phi_{pq},h_{pq})$ 满足下列条件，就称它为一个
\emph{坐标变换}：
\begin{description}[style=multiline,leftmargin=7.2em,labelwidth=6.5em]
\item[\textup{(A1.4.1)}] $h_{pq}\colon\Gamma_q\to\Gamma_p$ 是单射同态；
\item[\textup{(A1.4.2)}] $\phi_{pq}\colon V_{pq}\to V_p$ 是 $h_{pq}$-等变
光滑嵌入，其中 $V_{pq}$ 是 $o_q$ 的一个 $\Gamma_q$-不变开邻域；并且
诱导映射
$\underline\phi_{pq}\colon V_{pq}/\Gamma_q\to V_p/\Gamma_p$ 是单射；
\item[\textup{(A1.4.3)}] $(\widehat\phi_{pq},\phi_{pq})$ 是向量丛的
$h_{pq}$-等变嵌入
\[
 E_q\times V_{pq}\longrightarrow E_p\times V_p;
\]
\item[\textup{(A1.4.4)}]
$\widehat\phi_{pq}\circ s_q=s_p\circ\phi_{pq}$。这里以及下文中，
我们有时把 $s_p$ 看作平凡丛 $E_p\times V_p\to V_p$ 的截面
$s_p\colon V_p\to E_p\times V_p$；
\item[\textup{(A1.4.5)}] 在
$(s_q^{-1}(0)\cap V_{pq})/\Gamma_q$ 上，
$\psi_q=\psi_p\circ\underline\phi_{pq}$；这里的
$\underline\phi_{pq}$ 如 (A1.4.2)；
\item[\textup{(A1.4.6)}] 对任意 $x\in V_{pq}$，$h_{pq}$ 限制为同构
\[
 (\Gamma_q)_x\longrightarrow(\Gamma_p)_{\phi_{pq}(x)}.
\]
\end{description}
\end{definition}

结合条件 (A1.4.6) 与
$\underline\phi_{pq}\colon V_{pq}/\Gamma_q\to V_p/\Gamma_p$ 的单射性，
可得
\begin{description}[style=multiline,leftmargin=7.2em,labelwidth=6.5em]
\item[\textup{(A1.4.6$'$)}]
若 $\gamma\in\Gamma_p$ 且
$(\gamma\phi_{pq}(V_{pq}))\cap\phi_{pq}(V_{pq})\ne\varnothing$，则
$\gamma\in h_{pq}(\Gamma_q)$。
\end{description}
若把 $V_q$ 取得充分小，则 (A1.4.6) 等价于要求 $h_{pq}$ 限制为同构
\[
 \Gamma_q=(\Gamma_q)_{o_q}\longrightarrow
 (\Gamma_p)_{\phi_{pq}(o_q)}.
\]
在我们对模空间构造仓西结构时，$\Gamma_p$ 对应于稳定映射（可以带标记点）
的自同构群。因此，在我们的情形中 (A1.4.6) 成立。

\begin{definition}{A1.5}\label{def:a1-5}
$X$ 上的一个\emph{仓西结构}\liftIndex{I037}，是给每个 $p\in X$ 指定
一个仓西邻域 $(V_p,E_p,\Gamma_p,\psi_p,s_p)$，并给每个
$q\in\psi_p(s_p^{-1}(0)/\Gamma_p)$ 指定一个坐标变换
$(\widehat\phi_{pq},\phi_{pq},h_{pq})$，使下列条件成立：
\begin{description}[style=multiline,leftmargin=7.2em,labelwidth=6.5em]
\item[\textup{(A1.6.1)}] $\dim V_p-\operatorname{rank}E_p$ 与 $p$ 无关；
\item[\textup{(A1.6.2)}] 若
\[
 r\in\psi_q\bigl((V_{pq}\cap s_q^{-1}(0))/\Gamma_q\bigr),
 \qquad q\in\psi_p(s_p^{-1}(0)/\Gamma_p),
\]
则存在 $\gamma_{pqr}\in\Gamma_p$，使
\[
\begin{aligned}
 h_{pq}\circ h_{qr}
 &=\gamma_{pqr}\cdot h_{pr}\cdot\gamma_{pqr}^{-1},\\
 \phi_{pq}\circ\phi_{qr}
 &=\gamma_{pqr}\cdot\phi_{pr},\\
 \widehat\phi_{pq}\circ\widehat\phi_{qr}
 &=\gamma_{pqr}\cdot\widehat\phi_{pr}.
\end{aligned}
\]
其中第二个等式在
$\phi_{qr}^{-1}(V_{pq})\cap V_{qr}\cap V_{pr}$ 上成立，第三个等式在
$E_r\times(\phi_{qr}^{-1}(V_{pq})\cap V_{qr}\cap V_{pr})$ 上成立。
\end{description}
若 $V_p$ 有边界或角，则称
$(V_p,E_p,\Gamma_p,\psi_p,s_p)$ 为带边界或角的仓西结构。
\end{definition}

我们指出，(A1.6.2) 等价于
\[
 \underline\phi_{pq}\circ\underline\phi_{qr}
 =\underline\phi_{pr}.
\]
（利用 $\Gamma_p$ 作用的有效性即可证明这一等价性。）

可以定义仓西结构的芽。不过，本书不使用这一概念。我们约定两个仓西结构
\[
\begin{split}
 &\bigl(\{(V_p,E_p,\Gamma_p,\psi_p,s_p)_p\mid p\in X\},\\
 &\hspace{4em}\{(\widehat\phi_{pq},\phi_{pq},h_{pq})
   \mid p,q\in\psi_p(s_p^{-1}(0)/\Gamma_p)\}\bigr)
\end{split}
\]
与
\[
\begin{split}
 &\bigl(\{(V'_p,E'_p,\Gamma'_p,\psi'_p,s'_p)_p\mid p\in X\},\\
 &\hspace{4em}\{(\widehat\phi'_{pq},\phi'_{pq},h'_{pq})
   \mid p,q\in\psi'_p((s'_p)^{-1}(0)/\Gamma'_p)\}\bigr)
\end{split}
\]
相同，当且仅当所有资料都相同，即
$V_p=V'_p,E_p=E'_p,\ldots,h_{pq}=h'_{pq}$。
我们把 $\dim V_p-\operatorname{rank}E_p$ 称为仓西结构的\emph{虚维数}
（或简称维数）。

设 $K\subset X$，且
$\mathcal U=\{V_{p_i}/\Gamma_{p_i}\}$ 是一族仓西邻域。若
\[
 K\subset\bigcup_i\psi_{p_i}(s_{p_i}^{-1}(0)/\Gamma_{p_i}),
\]
则滥用记号，有时也称 $\mathcal U$ 为 $K$ 的一个仓西邻域。

依定义，$X$ 上的一个\emph{轨形结构}，就是满足每个 $p\in X$ 都有
$E_p=0$ 的仓西结构。

\begin{example}{A1.7}\label{ex:a1-7}
设 $W_1$ 是巴拿赫流形，$W_2$ 是巴拿赫空间，并设
$F\colon W_1\to W_2$ 是光滑（非线性）弗雷德霍姆映射。假设离散群 $G$
作用在 $W_1$ 上，并在 $W_2$ 上有线性表示，使 $F$ 等变。

考虑 $X$ 是方程
\begin{equation}
 F(w)=0
 \tag{A1.8}\label{eq:a1-8}
\end{equation}
之解的模空间这一情形。置
\[
 \widetilde X=\{w\in W_1\mid F(w)=0\}.
\]
假设 $G$ 在式~\eqref{eq:a1-8} 的解集上的作用正常不连续。（注意，
式~\eqref{eq:a1-8} 的解集局部紧，因为 $F$ 是弗雷德霍姆映射，从而局部
恰当。）现在考虑商空间 $X=\widetilde X/G$。通常，式~\eqref{eq:a1-8}
是非线性椭圆偏微分方程。

在这种情形下，可以定义 $X$ 的仓西结构。设 $w$ 是式~\eqref{eq:a1-8}
的一个解。如下考虑它的线性化方程：对 $u\in T_wW_1$，置
\[
 D_wF(u)=\left.\frac{d}{dt}F(w+tu)\right|_{t=0}.
\]
这里 $D_wF\colon T_wW_1\to W_2$；我们把 $T_0W_2\cong W_2$ 识别，
并回忆 $W_2$ 被假设为巴拿赫空间。

若 $D_wF\colon T_wW_1\to W_2$ 满射，则由隐函数定理，
$\widetilde X$ 在 $w$ 的邻域中是光滑流形。因此，可以把 $[w]$ 的仓西邻域
$V_w$ 取作 $\widetilde X$ 中 $w$ 的一个邻域，并取 $E_w=0$。

现在考虑 $D_wF\colon T_wW_1\to W_2$ 不是满射、但为弗雷德霍姆算子的
情形。选取有限维线性子空间 $E_w\subset W_2$，使
\begin{equation}
 \operatorname{Im}(D_wF)+E_w=W_2.
 \tag{A1.9}\label{eq:a1-9}
\end{equation}
定义 $V_w$ 为 $w$ 的一个邻域 $U(w)$ 中满足
\begin{equation}
 F(u)\equiv0\pmod{E_w}
 \tag{A1.10}\label{eq:a1-10}
\end{equation}
的所有 $u$ 所成集合。$\Gamma_w$ 是 $w$ 的自同构所成集合，即
\[
 \Gamma_w=\{g\in G\mid gw=w\}.
\]
由于假设 $G$ 在式~\eqref{eq:a1-8} 的解集上的作用正常不连续，
$\Gamma_w$ 是有限群。若 $u$ 满足式~\eqref{eq:a1-10}，置
\[
 s_w(u)=F(u)\in E_w.
\]
于是 $s_w$ 定义映射 $V_w\to E_w$。若 $u\in V_w$ 且 $s_w(u)=0$，
则 $u$ 是式~\eqref{eq:a1-8} 的解，因而其等价类 $\psi_w(u)$ 是 $X$ 的
元素。可以把 $V_w$ 取得充分小，使式~\eqref{eq:a1-9} 与隐函数定理蕴含
$V_w$ 是光滑流形。这样便得到了仓西邻域
$(V_w,E_w,\Gamma_w,\psi_w,s_w)$。

若 $w'$ 位于 $\psi_w$ 的像中，则可以选取 $E_{w'}\subset E_w$（见
第~\ref{subsec:7-1-4} 小节末尾的讨论）。这时
$F(u)\equiv0\pmod{E_{w'}}$ 蕴含式~\eqref{eq:a1-10}，故
$V_{w'}\cap U(w)$ 是 $V_w$ 的子流形。置
$V_{ww'}=V_{w'}\cap U(w)$，便有
$\phi_{ww'}\colon V_{ww'}\to V_w$。又因 $E_{w'}\subset E_w$，存在丛嵌入
$\widehat\phi_{ww'}$。最后，必要时缩小 $V_w$，使 $w'\in V_w$ 的自同构群
$\Gamma_{w'}$ 包含于 $w$ 的自同构群 $\Gamma_w$。由此得到坐标变换
$(\widehat\phi_{ww'},\phi_{ww'},h_{ww'})$。
\end{example}

一般地，$\dim V_p\ne\dim V_q$，所以坐标变换可能只在一个方向上存在。
为了得到带仓西结构空间的基本链，我们的思路是扰动仓西映射 $s_p$，使它
横截于零。扰动后的仓西映射应与坐标变换相容。与通常的横截性定理证明
相同，我们将在各仓西邻域上归纳构造仓西映射的扰动。由于坐标变换可能
只在一个方向上存在，必须巧妙选择仓西邻域的次序，才能完成这一归纳。
下面的良好坐标系正提供了这样的选择。

\begin{lemma}{A1.11}\label{lem:a1-11}
（\textbf{[FuOn99II]} 引理 6.3）设 $X$ 是带仓西结构的空间。则存在有限集
$P\subset X$、$P$ 上的一个次序 $<$，以及每个 $p\in P$ 的仓西邻域
$(V_p,E_p,\Gamma_p,\psi_p,s_p)$，满足下列性质：

\begin{description}[style=multiline,leftmargin=8.2em,labelwidth=7.5em]
\item[\textup{(A1.12.1)}] 若 $q<p$ 且
\[
 \psi_p(s_p^{-1}(0)/\Gamma_p)
 \cap\psi_q(s_q^{-1}(0)/\Gamma_q)\ne\varnothing,
\]
则存在 $(V_{pq},\widehat\phi_{pq},\phi_{pq},h_{pq})$，其中：
\begin{description}[style=multiline,leftmargin=9.2em,labelwidth=8.5em]
\item[\textup{(A1.12.1.1)}] $V_{pq}$ 是 $V_q$ 的 $\Gamma_q$-不变开子集，且
$V_{pq}/\Gamma_q$ 包含
\[
 \psi_q^{-1}\bigl(
  \psi_p(s_p^{-1}(0)/\Gamma_p)
  \cap\psi_q(s_q^{-1}(0)/\Gamma_q)
 \bigr);
\]
\item[\textup{(A1.12.1.2)}] $h_{pq}\colon\Gamma_q\to\Gamma_p$ 是单射同态，
其像为 $(\Gamma_p)_{\phi_{pq}(o_q)}$；\translatornote{条件 (A1.12.1.2) 的源书下标为
$(\Gamma_p)_{\phi_{pq}(q)}$；但 $q\in X$ 不能代入定义在 $V_q$ 上的
$\phi_{pq}$。依定义 A1.1、A1.3，改按选定中心点 $o_q$ 理解。}
\item[\textup{(A1.12.1.3)}] $\phi_{pq}\colon V_{pq}\to V_p$ 是
$h_{pq}$-等变光滑嵌入，并且诱导映射
$V_{pq}/\Gamma_q\to V_p/\Gamma_p$ 是单射；
\item[\textup{(A1.12.1.4)}]
$(\widehat\phi_{pq},\phi_{pq})$ 是向量丛的 $h_{pq}$-等变嵌入
$E_q\times V_{pq}\to E_p\times V_p$；
\item[\textup{(A1.12.1.5)}]
$\widehat\phi_{pq}\circ s_q=s_p\circ\phi_{pq}$，且
$\psi_q=\psi_p\circ\underline\phi_{pq}$。
\end{description}

\item[\textup{(A1.12.2)}] 若 $r<q<p$ 且
\[
 \psi_p(s_p^{-1}(0)/\Gamma_p)
 \cap\psi_q(s_q^{-1}(0)/\Gamma_q)
 \cap\psi_r(s_r^{-1}(0)/\Gamma_r)\ne\varnothing,
\]
则存在 $\gamma_{pqr}\in\Gamma_p$，使
\[
\begin{aligned}
 h_{pq}\circ h_{qr}
 &=\gamma_{pqr}\cdot h_{pr}\cdot\gamma_{pqr}^{-1},\\
 \phi_{pq}\circ\phi_{qr}
 &=\gamma_{pqr}\cdot\phi_{pr},\\
 \widehat\phi_{pq}\circ\widehat\phi_{qr}
 &=\gamma_{pqr}\cdot\widehat\phi_{pr}.
\end{aligned}
\]
第二个等式在 $\phi_{qr}^{-1}(V_{pq})\cap V_{qr}\cap V_{pr}$ 上成立，
第三个等式在
$E_r\times(\phi_{qr}^{-1}(V_{pq})\cap V_{qr}\cap V_{pr})$ 上成立。

\item[\textup{(A1.12.3)}]
\[
 \bigcup_{p\in P}\psi_p(s_p^{-1}(0)/\Gamma_p)=X.
\]
\end{description}
\end{lemma}

我们把 $\{(V_p,E_p,\Gamma_p,\psi_p,s_p)\}_{p\in P}$ 称为一个
\emph{良好坐标系}\liftIndex{I309}。
引理~\ref{lem:a1-11} 的证明颇具技术性，这里从略；见
\textbf{[FuOn99II]} 第 957--958 页。

虚基本链可以看作某个无穷维空间中的链，而所有仓西邻域都嵌入该空间。
不过，本书的目的并不需要这一图景。我们使用空间 $X$ 上的仓西结构，并有
一个从 $X$ 到另一空间（记作 $Y$）的映射；我们把虚基本链用作 $Y$ 中的链。
在本书的语境中，这一映射将是从模空间出发的评价映射。具体地说，我们
在以下情形中工作。

\begin{definition}{A1.13}\label{def:a1-13}
考虑引理~\ref{lem:a1-11} 的情形，并设 $Y$ 是拓扑空间。若一族
$\Gamma_p$-等变连续映射 $f_p\colon V_p\to Y$ 满足
\[
 f_p\circ\phi_{pq}=f_q
\]
于 $V_{pq}$ 上成立，就称 $\{f_p\}$ 为一个\emph{强连续映射}\liftIndex{I372}。
强连续映射诱导连续映射 $f\colon X\to Y$。意义清楚时，我们也含混地把
$f=\{f_p\}$ 记作同一个符号。

若 $Y$ 是光滑流形，则称强连续映射 $f\colon X\to Y$ \emph{光滑}，
如果每个 $f_p\colon V_p\to Y$ 都光滑；称它\emph{弱浸没}\liftIndex{I381}，如果每个
$f_p$ 都是浸没。
\end{definition}

在本书所用构造中，从各个模空间出发的评价映射都是弱浸没的。特别地，
评价映射在模空间每个边界或角上的限制也都是弱浸没的。

为了使带仓西结构空间 $X$ 承载基本链，还需要若干附加条件。设
\[
 x\in
 \psi_p(s_p^{-1}(0)/\Gamma_p)
 \cap\psi_q(s_q^{-1}(0)/\Gamma_q)\subset X,
\]
其中 $(V_p,E_p,\Gamma_p,\psi_p,s_p)$ 与
$(V_q,E_q,\Gamma_q,\psi_q,s_q)$ 是仓西邻域。假设
$x=\psi_p(x_p)=\psi_q(x_q)$。考虑
\[
 d_{x_p}s_p\colon T_{x_p}V_p\longrightarrow(E_p)_{x_p},
 \qquad
 d_{x_q}s_q\colon T_{x_q}V_q\longrightarrow(E_q)_{x_q}.
\]
可以把它们看作我们的“模空间”的切复形。事实上，若 $s_p$ 在 $x_p$
处横截于 $0$，则 $X$ 在 $x_p$ 的一个邻域中是流形，并且
$T_{x_p}X=\ker(d_{x_p}s_p)$。下面的条件蕴含：切复形
$d_{x_p}s_p\colon T_{x_p}V_p\to(E_p)_{x_p}$ 可以胶合成 $X$ 上的一个
虚向量丛。

考虑引理~\ref{lem:a1-11} 的情形。\translatornote{源书此处写“引理 A1.10 的情形”，但本书没有引理 A1.10；
(A1.10) 是例 A1.7 的同余方程，而这里使用的是引理 A1.11 的坐标变换资料。}
利用适当黎曼度量的指数映射，把
$V_p$ 中 $\phi_{pq}(V_{pq})$ 的一个邻域，与法丛
$NV_{pq}V_p\to V_{pq}$ 的零截面的一个邻域识别。沿纤维方向对仓西映射
$s_p$ 取纤维导数，得到 $h_{pq}$-等变丛同态
\[
 d_{\mathrm{fiber}}s_p\colon NV_{pq}V_p
 \longrightarrow E_p\times V_{pq}.
\]

\begin{definition}{A1.14}\label{def:a1-14}
若 $d_{\mathrm{fiber}}s_p$ 诱导 $V_{pq}$ 上 $\Gamma_q$-等变向量丛同构
\begin{equation}
 NV_{pq}V_p\cong
 \frac{E_p\times V_{pq}}
      {\widehat\phi_{pq}(E_q\times V_{pq})},
 \tag{A1.15}\label{eq:a1-15}
\end{equation}
就称带仓西结构空间 $X$ \emph{有切丛}\liftIndex{I312}。
\end{definition}

注意，条件 (A1.4.6) 保证光滑轨形嵌入
$V_{pq}/\Gamma_q\to V_p/\Gamma_p$ 的法丛可定义为轨形丛（轨形丛的定义
见定义~\ref{def:a1-33}）。我们还指出，同构~\eqref{eq:a1-15} 不依赖于
把给定管状邻域与法丛开集识别时所选的黎曼度量。

依定义，对每个
$x\in\phi_{qr}^{-1}(V_{pq})\cap V_{qr}\cap V_{pr}$，下图交换：
\begin{equation}
\resizebox{0.96\textwidth}{!}{$
\begin{array}{ccccc}
 (NV_{qr}V_q)_x
 &\xrightarrow{\ \gamma_{pqr}\cdot d\phi_{pq}\ }
 &(NV_{pr}V_p)_x
 &\longrightarrow
 &(NV_{pq}V_p)_{\phi_{qr}(x)}
 \\[0.7em]
 \Big\downarrow &&\Big\downarrow&&\Big\downarrow
 \\[0.7em]
 \displaystyle\frac{E_q}{\widehat\phi_{qr,x}(E_r)}
 &\xrightarrow{\ \gamma_{pqr}\cdot\widehat\phi_{pq}\ }
 &\displaystyle\frac{E_p}{\widehat\phi_{pr,x}(E_r)}
 &\longrightarrow
 &\displaystyle\frac{E_p}{\widehat\phi_{pq,\phi_{qr}(x)}(E_q)}
\end{array}$}
\tag{A1.16}\label{diag:a1-16}
\end{equation}
这里以及下文中，
$\widehat\phi_{pq,x}\colon E_q\to E_p$ 表示丛映射
$\widehat\phi_{pq}$ 在 $x$ 的纤维上的限制。

\begin{definition}{A1.17}\label{def:a1-17}
设 $X$ 是有切丛的带仓西结构空间。若有
\[
 \Lambda^{\mathrm{top}}E_p^*\otimes\Lambda^{\mathrm{top}}TV_p
\]
的一个平凡化，它与同构~\eqref{eq:a1-15} 相容，并且 $\Gamma_p$ 的作用
保持该平凡化的同伦类，就称 $X$ 上的仓西结构\emph{已定向}%
\liftIndex{I311}。
\end{definition}

\begin{example}{A1.18}\label{ex:a1-18}
考虑例~\ref{ex:a1-7} 的情形。我们说明这一仓西结构有切丛。设 $w'$ 位于
$\psi_w$ 的像中。已有 $E_{w'}\subset E_w$。因此，若
$u\in V_{ww'}\subset V_w$，则
\[
 T_uV_w=(D_uF)^{-1}(E_w),
 \qquad
 T_uV_{w'}=(D_uF)^{-1}(E_{w'}).
\]
于是，由于 $\operatorname{Im}(D_uF)+E_{w'}=W_2$，$dF$ 诱导同构
\[
 \frac{T_uV_w}{T_uV_{w'}}\cong\frac{E_w}{E_{w'}}.
\]
同构~\eqref{eq:a1-15} 由这一同构诱导。本例中的仓西映射是 $F$ 的限制，
所以以上同构也由仓西映射诱导\translatornote{源书在本例的两条切空间公式中写 $D_uP$，随后又写
$\operatorname{Im}(D_vP)$；但例 A1.7 唯一给定的映射是 $F$，本段末句也
明确写 $dF$。故三处均按 $D_uF$ 排印。}。
\end{example}

现在说明 \textbf{[FuOn99II]} 中虚基本链的构造。为此需要多值截面的概念。
\liftIndexRangeStart{I314}\liftIndexRangeStart{I315}%
假设 $\Gamma$ 作用在流形 $V$ 与向量空间 $E$ 上。阶为 $n!$ 的对称群
$\mathfrak S_n$ 按下式作用在乘积 $E^n$ 上：
\[
 \sigma(x_1,\ldots,x_n)
 =(x_{\sigma(1)},\ldots,x_{\sigma(n)}).
\]
令 $S^n(E)$ 为商空间 $E^n/\mathfrak S_n$。于是，$\Gamma$ 在 $E$ 上的
作用诱导 $S^n(E)$ 上的作用。映射 $E^n\to E^{nm}$
\[
 (x_1,\ldots,x_n)\longmapsto
 \bigl(
 \underbrace{x_1,\ldots,x_1}_{m\text{ 次}},
 \underbrace{x_2,\ldots,x_2}_{m\text{ 次}},\ldots,
 \underbrace{x_n,\ldots,x_n}_{m\text{ 次}}
 \bigr)
\]
诱导 $\Gamma$-等变映射 $S^n(E)\to S^{nm}(E)$。

\begin{definition}{A1.19}\label{def:a1-19}
平凡丛 $\pi\colon E\times V\to V$ 的一个 $n$-\emph{多值截面} $s$，
是 $\Gamma$-等变映射 $V\to S^n(E)$。若存在
\[
 \widetilde s=(\widetilde s_1,\ldots,\widetilde s_n)\colon V\to E^n,
\]
使它与商投影 $E^n\to S^n(E)$ 的复合等于 $s$，就称 $s$ \emph{可提升}。
（不要求 $\widetilde s$ 为 $\Gamma$-等变。）每个
$\widetilde s_1,\ldots,\widetilde s_n$ 都称为 $s$ 的一个\emph{分支}。

若 $s\colon V\to S^n(E)$ 是 $n$-多值截面，则对每个 $m$，把它与
$S^n(E)\to S^{nm}(E)$ 复合，便得到一个 $nm$-多值截面。

若 $n$-多值截面 $s$ 与 $m$-多值截面 $s'$ 所诱导的两个 $nm$-多值截面
相同，就称 $s$ 与 $s'$ 等价。这个等价关系下的等价类称为一个
\emph{多值截面}。

若可提升多值截面的每个分支都横截于零，就称该多值截面横截于零。

称一族多值截面 $s_\epsilon$ 在 $\epsilon\to0$ 时收敛到 $s$，如果存在
$n$，使每个 $s_\epsilon$ 都可由一个 $n$-多值截面
$s_\epsilon^n$ 表示，并且 $s_\epsilon^n$ 收敛到 $s$ 的一个代表。
\end{definition}

除非另有说明，从现在起假定所有多值截面都可提升。

\begin{lemma}{A1.20}\label{lem:a1-20}
设平凡丛 $\pi\colon E\times V\to V$ 的多值截面 $s$ 在
$K\subset V/\Gamma$ 的一个仓西邻域中横截于零。则存在一族多值截面
$s_\epsilon$，它在 $C^0$ 意义下收敛到 $s$，每个 $s_\epsilon$ 都横截于零，
并且在 $K$ 的一个邻域中 $s_\epsilon=s$。
\end{lemma}
\begin{proof}
令 $m$ 为 $\Gamma$ 的阶。设 $s$ 由
$E^n\to S^n(E)$ 与 $\widetilde s^{\,n}\colon V\to E^n$ 的复合表示。
由通常的横截性定理，存在
\[
 \widetilde s_\epsilon
 =(\widetilde s_{\epsilon,i})_{i=1,\ldots,n},
\]
它收敛到 $\widetilde s^{\,n}$，横截于零，并在 $\widetilde K$ 上等于
$\widetilde s^{\,n}$；这里 $\widetilde K\subset V$ 且
$\widetilde K/\Gamma=K$。定义 $\widetilde s'_\epsilon\colon V\to E^{nm}$ 为
\[
 \widetilde s'_\epsilon
 =\bigl(
   \sigma\circ\widetilde s_{\epsilon,i}\circ\sigma^{-1}
  \bigr)_{i=1,\ldots,n,\,\sigma\in\Gamma}.
\]
（这里 $m=\#\Gamma$。）$\widetilde s'_\epsilon$ 的等价类就是所求多值截面。
\end{proof}

\begin{definition}{A1.21}\label{def:a1-21}
考虑引理~\ref{lem:a1-11} 的情形，并假设仓西结构有切丛。设
$p,q\in P$、$q<p$，并分别给定 $E_p\times V_p$ 与 $E_q\times V_q$ 的
多值截面 $s'_p,s'_q$。下面定义二者的相容性。

考虑 (A1.12.1.3) 中的嵌入 $\phi_{pq}\colon V_{pq}\to V_p$，把它的法丛
$NV_{pq}V_p$ 与 $\phi_{pq}(V_{pq})$ 的一个管状邻域识别。对每个
$x\in V_{pq}$ 固定分裂
\begin{equation}
 E_p\cong
 \widehat\phi_{pq,x}(E_q)
 \oplus\frac{E_p}{\widehat\phi_{pq,x}(E_q)}.
 \tag{A1.22}\label{eq:a1-22}
\end{equation}
对每个 $y\in NV_{pq}V_p$，利用同构~\eqref{eq:a1-15} 得到一个元素
\[
 1(y)\in\frac{E_p}{\widehat\phi_{pq,\pi(y)}(E_q)},
\]
其中 $\pi\colon NV_{pq}V_p\to V_{pq}$ 是投影\translatornote{源书在此商纤维中写 $\pi(x)$；当前量化的法向量是 $y$，
而紧随公式也使用 $\pi(y)$，故按 $\pi(y)$ 排印。}。

现在如下在 $NV_{pq}V_p$ 上定义多值截面 $s'_q\oplus1$。设
$(\widetilde s'_{q,1},\ldots,\widetilde s'_{q,n})$ 是 $s'_q$ 的一个代表。
利用式~\eqref{eq:a1-22}，置
\[
\begin{aligned}
 (\widetilde s'_q\oplus1)(y)
 &=\bigl(
   \widetilde s'_{q,1}(\pi(y))\oplus1(y),\ldots,
   \widetilde s'_{q,n}(\pi(y))\oplus1(y)
  \bigr)\\
 &\in
 \left(
  \widehat\phi_{pq,\pi(y)}(E_q)
  \oplus\frac{E_p}{\widehat\phi_{pq,\pi(y)}(E_q)}
 \right)^n
 \cong E_p^n.
\end{aligned}
\]
以 $s'_q\oplus1$ 表示 $\widetilde s'_q\oplus1$ 的等价类。

若 $s'_p$ 在 $NV_{pq}V_p$ 上的限制等于 $s'_q\oplus1$，就称 $s'_p$ 与
$s'_q$ 相容。注意，单值截面 $s_p$（即仓西映射）按此定义与 $s_q$ 相容；
这是因为同构~\eqref{eq:a1-15} 由仓西映射的纤维导数诱导。

一个多值截面 $\mathfrak s=\{s'_p\}_{p\in P}$，是相容的多值截面系
（$s'_p\colon V_p\to E_p$）。若 $X$ 是带仓西结构空间，则依定义，
$X$ 的多值截面就是其障碍丛的多值截面。
\end{definition}
\liftIndexRangeEnd{I315}\liftIndexRangeEnd{I314}

\begin{theorem}{A1.23}\label{thm:a1-23}
考虑引理~\ref{lem:a1-11} 的情形。假设 $X$ 上的仓西结构有切丛。
则存在一族多值截面
$\mathfrak s_\epsilon=\{s'_{p,\epsilon}\}$，它收敛到
$\{s_p\}_{p\in P}$（即仓西映射），并且对所有 $\epsilon>0$，每个
$s'_{p,\epsilon}$ 都横截于 $0$。

进一步，若 $K\subset X$ 是紧集，而
$\{s''_{p,\epsilon}\}$ 是 $K$ 的一个仓西邻域中的一族多值截面，
它收敛到 $\{s_p\}_{p\in P}$ 且在 $K$ 上横截于 $0$，则可以选取
$\{s'_{p,\epsilon}\}_{p,\epsilon}$，使它在 $K$ 的某个仓西邻域上等于
$\{s''_{p,\epsilon}\}$。
\end{theorem}
\begin{proof}
按照 $P$ 上的次序 $<$，在各 $V_p$ 上归纳构造
$\{s'_{p,\epsilon}\}_{p,\epsilon}$。每个归纳步骤如下应用
引理~\ref{lem:a1-20}。假设 $q<p$ 时 $\{s'_{q,\epsilon}\}$ 已经确定。
那么在 $NV_{pq}V_p$（$q<p$）的并集上，多值截面
$s'_{p,\epsilon}$ 已由定义~\ref{def:a1-21} 所述的
$\widetilde s'_{q,\epsilon}\oplus1$ 确定。（若 $r<q<p$，则由
图式~\eqref{diag:a1-16} 的交换性，在
$NV_{pq}V_p\cap NV_{pr}V_p$ 上有
$\widetilde s'_{q,\epsilon}\oplus1
=\widetilde s'_{r,\epsilon}\oplus1$。）现在即可由引理~\ref{lem:a1-20}
构造 $\widetilde s'_{p,\epsilon}$。
\end{proof}

设 $Y$ 是拓扑空间，且 $\{f_p\}$（$f_p\colon V_p\to Y$）是强连续映射。
我们将使用定理~\ref{thm:a1-23} 所构造的多值截面
$\mathfrak s_\epsilon$，构造 $Y$ 中的一条链，即虚基本链。具体构造如下。
置
\[
 \widetilde{(s'_{p,\epsilon})^{-1}(0)}
 =\left\{y\in V_p\ \middle|\
 \begin{array}{c}
  s'_{p,\epsilon}\text{ 有一个分支 }
  \widetilde s'_{p,\epsilon,i},\\[-0.2em]
  \widetilde s'_{p,\epsilon,i}(y)=0
 \end{array}
 \right\},
\]
并令
\[
 (s'_{p,\epsilon})^{-1}(0)
 =\widetilde{(s'_{p,\epsilon})^{-1}(0)}/\Gamma_p.
\]
相容性蕴含
\begin{equation}
 \underline\phi_{pq}\left(
  \frac{
   \widetilde{(s'_{q,\epsilon})^{-1}(0)}\cap V_{pq}}
       {\Gamma_q}
 \right)
 =
 \frac{
  \widetilde{(s'_{p,\epsilon})^{-1}(0)}
  \cap\Gamma_p\cdot(NV_{pq}V_p)}
 {\Gamma_p},
 \tag{A1.24}\label{eq:a1-24}
\end{equation}
其中
$\underline\phi_{pq}\colon V_{pq}/\Gamma_q\to V_p/\Gamma_p$。置
\[
 \bigl((s'_\epsilon)^{-1}(0)\bigr)_{\mathrm{set}}
 =\left(
   \bigcup_{p\in P}(s'_{p,\epsilon})^{-1}(0)
  \right)\big/\sim,
\]
其中 $\sim$ 表示把
$[y]\in(s'_{q,\epsilon})^{-1}(0)$ 与
$\underline\phi_{pq}([y])\in(s'_{p,\epsilon})^{-1}(0)$ 识别\translatornote{式 (A1.24) 的源书右端写
$\Gamma_p\cdot(NV_{pq}V_q)$；但该项须位于 $V_p$ 并与
$(s'_{p,\epsilon})^{-1}(0)$ 相交，且定义 A1.21 始终使用
$NV_{pq}V_p$，故恢复其环境空间为 $V_p$。}。

对每个 $p$，选取 $n_p$ 以及
\[
 \widetilde s'_{p,\epsilon}
 =(\widetilde s'_{p,\epsilon,1},\ldots,
   \widetilde s'_{p,\epsilon,n_p})
 \colon V_p\longrightarrow E_p^{n_p},
\]
使 $\widetilde s'_{p,\epsilon}$ 与投影
$E_p^{n_p}\to S^{n_p}(E_p)$ 的复合代表 $s'_{p,\epsilon}$。对
$y\in\widetilde{(s'_{p,\epsilon})^{-1}(0)}$，置
\begin{equation}
 \operatorname{val}_p(y)
 =\#\{i\mid\widetilde s'_{p,\epsilon,i}(y)=0\}.
 \tag{A1.25}\label{eq:a1-25}
\end{equation}
它在 $\Gamma_p$ 作用下不变，因而诱导
\[
 \operatorname{val}_p\colon
 (s'_{p,\epsilon})^{-1}(0)\longrightarrow\mathbb Z_{>0}.
\]

\begin{lemma}{A1.26}\label{lem:a1-26}
（\textbf{[FuOn99II]} 引理 6.9）对 $s'_\epsilon$ 的一般选择，空间
$\bigl((s'_\epsilon)^{-1}(0)\bigr)_{\mathrm{set}}$ 有一个光滑三角剖分，
使得对剖分中的每个单形 $\Delta_a^m$，都存在 $p_a$，满足
\[
 \Delta_a^m\subset(s'_{p_a,\epsilon})^{-1}(0),
\]
并且 $\operatorname{val}_{p_a}$ 在 $\operatorname{Int}\Delta_a^m$ 上为常数。
\end{lemma}

证明颇具技术性，这里从略；读者可参见 \textbf{[FuOn99II]} 第 946 页。

现在假设 $X$ 已定向，并令 $d$ 为 $X$ 的虚维数。对
$\bigl((s'_\epsilon)^{-1}(0)\bigr)_{\mathrm{set}}$ 的每个 $d$ 维单形
$\Delta_a^d$，如下定义其重数。（注意，由横截性，
$\bigl((s'_\epsilon)^{-1}(0)\bigr)_{\mathrm{set}}$ 的维数为 $d$。）

\begin{definition}{A1.27}\label{def:a1-27}
考虑引理~\ref{lem:a1-26} 的情形。设
\[
 [y]\in\operatorname{Int}\Delta_a^d
 \subset V_{p_a}/\Gamma_{p_a},
 \qquad y\in V_{p_a}.
\]
令
\[
 \widetilde s'_{p_a,\epsilon}
 =(\widetilde s'_{p_a,\epsilon,1},\ldots,
   \widetilde s'_{p_a,\epsilon,n_{p_a}})
 \colon V_{p_a}\longrightarrow E_{p_a}^{n_{p_a}}
\]
如 $\operatorname{val}_{p_a}$ 的定义。对每个满足
$\widetilde s'_{p_a,\epsilon,i}(y)=0$ 的 $i$，如下定义
$\epsilon_i\in\{\pm1\}$。由于 $s'_\epsilon$ 横截于 $0$，映射
\[
 d_y\widetilde s'_{p_a,\epsilon,i}\colon
 T_yV_{p_a}\longrightarrow E_{p_a}
\]
满射。利用定向所诱导的
\[
 \Lambda^{\mathrm{top}}E_{p_a}^*
 \otimes\Lambda^{\mathrm{top}}T_yV_{p_a}
\]
之平凡化，可以按照第 8 章约定~\ref{conv:8-2-1}(4)，给
$(d_y\widetilde s'_{p_a,\epsilon,i})^{-1}(0)$ 定向。若这一取向与
$\Delta_a^d$ 的取向相同，就令 $\epsilon_i=1$；否则令
$\epsilon_i=-1$。定义
\[
 \operatorname{mul}_{\Delta_a^d}
 =\frac{\displaystyle\sum_i\epsilon_i}{n_{p_a}}
 \in\mathbb Q,
\]
其中分子中的和取遍所有满足
$\widetilde s'_{p_a,\epsilon,i}(y)=0$ 的 $i$。

由 $\operatorname{val}_{p_a}$ 为常数，可以证明
$\operatorname{mul}_{\Delta_a^d}$ 与
$[y]\in\Delta_a^d$ 的选择无关。由相容性，可以证明它也与满足
$\Delta_a^d\subset(s'_{p_a,\epsilon})^{-1}(0)$ 的 $p_a$ 的选择无关\translatornote{源书在本定义开头把商群写成 $\Gamma_p$，末句又把固定的
$d$ 维单形写成 $\Delta_a^m$；依同句的 $V_{p_a}$、
$\operatorname{val}_{p_a}$ 及全段维数记号，分别恢复为
$\Gamma_{p_a}$ 与 $\Delta_a^d$。}。
\end{definition}

现在定义 $Y$ 中的虚基本链 $f_*[X]$。

\begin{definition}{A1.28}\label{def:a1-28}
\liftIndex{I105}\liftIndex{I313}%
\begin{equation}
 f_*[X]=\sum_a
 \operatorname{mul}_{\Delta_a^d},(f_{p_a})_*[\Delta_a^d].
 \tag{A1.29}\label{eq:a1-29}
\end{equation}
\end{definition}
这里 $f_*[X]$ 是 $Y$ 的一条 $\mathbb Q$-奇异链。

接着考虑 $X$ 有边界或角的情形。先回顾 $n$ 维带角流形 $V$。令
$S_kV$ 为 $V$ 的余维 $k$ 角中各点所成集合，$\overline S_kV$ 为其闭包。
换言之，若 $x$ 有坐标卡 $[0,1)^k\times\mathbb R^{n-k}$，且 $x$ 被映到
$0$，则 $x\in S_kV$。存在一个 $n-k$ 维带角流形
$\widehat S_kV$ 与连续满射
\[
 \pi\colon\widehat S_kV\longrightarrow\overline S_kV,
\]
使 $\pi$ 整体上有限对一，而在 $S_kV$ 上一对一。注意，第~\ref{subsec:7-1-1}
小节的几何情形中也出现过这样的有限对一映射。从现在起，我们滥用记号，
把 $\widehat S_kV$ 与 $\overline S_kV$ 识别。为了定义虚基本链，无须区分
二者，因为它们在顶层 $S_kV$ 上同构。下一定义所述仓西结构的情形也作
同样约定。

\begin{definition}{A1.30}\label{def:a1-30}
若 $x=\psi_p(\pi\widetilde x)$，其中 $\widetilde x\in V_p$ 位于 $V_p$ 的
余维 $k$ 角，且 $\pi\colon V_p\to V_p/\Gamma_p$ 是投影，就称
$x\in X$ 位于 $X$ 的余维 $k$ 角。（注意，$V_p$ 是带角流形。）以 $S_kX$
表示 $X$ 的余维 $k$ 角中所有元素所成集合，以 $\overline S_kX$ 表示其
闭包。记
\[
 \partial X=\overline S_1X,
\]
即 $X$ 的边界。

如下在 $\overline S_kX$ 上定义仓西结构。若
$p\in\overline S_kX$，而 $\Gamma_p$ 的作用在
$\overline S_kV_p$ 上诱导有效作用，就取
\[
 (\overline S_kV_p,E_p,\Gamma_p,
   s_p|_{\overline S_kV_p})
\]
为其仓西邻域。

若 $\Gamma_p$ 在 $\overline S_kV_p$ 上的诱导作用并不有效，则同时给
仓西邻域与障碍丛加上 $\Gamma_p$ 的同一个表示，使它在仓西邻域上的作用
变为有效。
\end{definition}

\begin{lemma}{A1.31}\label{lem:a1-31}
有
\[
 \partial f_*[X]=(f|_{\partial X})_*[\partial X].
\]
\end{lemma}
证明很容易，留给读者。由引理~\ref{lem:a1-31} 可证明如下结论。

\begin{lemma}{A1.32}\label{lem:a1-32}
若 $\partial X=\varnothing$，则 $f_*[X]$ 是闭链。它的同调类与扰动
$\mathfrak s_\epsilon$ 的选择无关。
\end{lemma}
\begin{proof}
前半立即由引理~\ref{lem:a1-31} 得出。为证后半，注意
$X\times[0,1]$ 有带边界的仓西结构，并且
\[
 \partial(X\times[0,1])=X\times\{0,1\}.
\]
假设 $\mathfrak s_\epsilon$ 有两种选择，记作
$\mathfrak s^0_\epsilon,\mathfrak s^1_\epsilon$。可以在
$X\times[0,1]$ 上找到 $\mathfrak s^+_\epsilon$，使它在
$X\times\{t_0\}$（$t_0=0,1$）上的限制等于
$\mathfrak s^{t_0}_\epsilon$（定理~\ref{thm:a1-23}）。此时后半也立即由
引理~\ref{lem:a1-31} 得出。
\end{proof}

下面回顾 \textbf{[FuOn99II]} 中丛系的定义。

\begin{definition}{A1.33}\label{def:a1-33}
设 $X$ 是带仓西结构空间。$X$ 上的一个\emph{丛系}由以下资料组成。
设 $(V_p,E_p,\Gamma_p,\psi_p,s_p)$ 是 $p\in X$ 的仓西邻域。
\begin{description}[style=multiline,leftmargin=7.2em,labelwidth=6.5em]
\item[\textup{(A1.34.1)}]
对每个 $p$，有 $V_p$ 上的一对向量丛 $F_{1,p},F_{2,p}$，而 $\Gamma_p$
作用在 $F_{1,p},F_{2,p}$ 上。

\item[\textup{(A1.34.2)}]
若 $q\in\psi_p(V_p/\Gamma_p)$，则对 $i=1,2$，有覆盖
$\phi_{pq}\colon V_q\to V_p$ 的 $h_{pq}$-等变向量丛嵌入
\[
 \widehat\phi_{i,pq}\colon F_{i,q}\longrightarrow F_{i,p}.
\]
此外有同构
\[
 \Phi_{pq}\colon
 \frac{F_{1,p}|_{\phi_{pq}(V_q)}}
      {\widehat\phi_{1,pq}(F_{1,q})}
 \xrightarrow{\ \cong\ }
 \frac{F_{2,p}|_{\phi_{pq}(V_q)}}
      {\widehat\phi_{2,pq}(F_{2,q})}.
\]

\item[\textup{(A1.34.3)}]
若
$r\in\psi_q(V_q/\Gamma_q)\subseteq\psi_p(V_p/\Gamma_p)$，则
\[
 \widehat\phi_{i,pq}\circ\widehat\phi_{i,qr}
 =\gamma_{pqr}\cdot\widehat\phi_{i,pr}.
\]

\item[\textup{(A1.34.4)}]
若
$r\in\psi_q(V_q/\Gamma_q)\subseteq\psi_p(V_p/\Gamma_p)$，则下图交换：\translatornote{条件 (A1.34.2) 的源书指标写作 $i=0,1$，但所定义的丛为
$F_{1,p},F_{2,p}$；条件 (A1.34.4) 的外层图表误写
$\psi_q(V_p/\Gamma_p)$；交换图第一列的两个分母又把取值于 $F_{i,q}$
的嵌入误写成 $\widehat\phi_{i,pr}$。分别恢复为 $i=1,2$、
$\psi_p(V_p/\Gamma_p)$ 与 $\widehat\phi_{i,qr}$。}
\end{description}
\[
\resizebox{0.98\textwidth}{!}{$
\begin{array}{ccccccccc}
0&\longrightarrow&
\displaystyle
\frac{F_{1,q}|_{\phi_{qr}(V_r)}}
     {\widehat\phi_{1,qr}(F_{1,r})}
&\xrightarrow{\widehat\phi_{1,pq}}&
\displaystyle
\frac{F_{1,p}|_{\gamma_{pqr}\cdot\phi_{pr}(V_r)}}
     {\gamma_{pqr}\cdot\widehat\phi_{1,pr}(F_{1,r})}
&\xrightarrow{\pi\circ\gamma_{pqr}^{-1}}&
\displaystyle
\frac{F_{1,p}|_{\phi_{pq}(V_q)}}
     {\widehat\phi_{1,pq}(F_{1,q})}
&\longrightarrow&0
\\[1.2em]
&&\underset{\Phi_{qr}}{\Big\downarrow}
&&\underset{\gamma_{pqr}\circ\Phi_{pr}\circ\gamma_{pqr}^{-1}}
             {\Big\downarrow}
&&\underset{\Phi_{pq}}{\Big\downarrow}&&
\\[1.2em]
0&\longrightarrow&
\displaystyle
\frac{F_{2,q}|_{\phi_{qr}(V_r)}}
     {\widehat\phi_{2,qr}(F_{2,r})}
&\xrightarrow{\widehat\phi_{2,pq}}&
\displaystyle
\frac{F_{2,p}|_{\gamma_{pqr}\cdot\phi_{pr}(V_r)}}
     {\gamma_{pqr}\cdot\widehat\phi_{2,pr}(F_{2,r})}
&\xrightarrow{\pi\circ\gamma_{pqr}^{-1}}&
\displaystyle
\frac{F_{2,p}|_{\phi_{pq}(V_q)}}
     {\widehat\phi_{2,pq}(F_{2,q})}
&\longrightarrow&0.
\end{array}$}
\]
\end{definition}

若 $X$ 是轨形，即对所有 $p$ 都有 $E_p=0$，而一对
$(F_{1,p},\widehat\phi_{1,pq})$ 在 $F_{2,p}=0$ 时满足 (A1.34)，就称它为
$X$ 上的一个\emph{轨形丛}。换言之，轨形丛是一对
$(\{F_{1,p}\},\{\widehat\phi_{1,pq}\})$，其中 $F_{1,p}$ 是 $V_p$ 上的
$\Gamma_p$-等变向量丛，$\widehat\phi_{1,pq}$ 是向量丛的
$h_{pq}$-等变同构，并且
\[
 \widehat\phi_{1,pq}\circ\widehat\phi_{1,qr}
 =\gamma_{pqr}\cdot\widehat\phi_{1,pr}.
\]

\begin{example}{A1.35}\label{ex:a1-35}
若 $X$ 上的仓西结构有切丛，则
\[
 F_{1,p}=TV_p,
 \qquad
 F_{2,p}=E_p\times V_p
\]
定义一个丛系。事实上，$\Phi_{pq}$ 正是同构~\eqref{eq:a1-15}；
(A1.34.4) 的交换性就是图式~\eqref{diag:a1-16} 的交换性。
\end{example}
